\section{The arithmetic foundation and its exact meaning}\label{sec:foundation}

\subsection{The base theory}
Throughout, \ipra\ means first-order intuitionistic primitive recursive arithmetic.
Its only variables range over natural numbers. Its terms contain $0$, successor,
and a symbol for each fixed primitive recursive (PR) definition; symbols come
with their defining equations. Its nonlogical axioms are these equations,
successor injectivity, $Sx\ne0$, and the universal closures of
\[
 [A(0,\vec u)\wedge\forall n(A(n,\vec u)\to A(Sn,\vec u))]
       \longrightarrow\forall n A(n,\vec u),                 \tag{I}
\]
where $A$ is quantifier-free. Logical rules are the usual intuitionistic
introduction and elimination rules for $\bot,\wedge,\vee,\to,\forall,\exists$,
with equality reflexivity and substitution. Negation abbreviates implication
to $\bot$. Quantifiers in a theorem need not be bounded or quantifier-free;
only the induction formula in (I) is restricted. This is also the base theory
specified by ProveIt's current iPRA sources \cite{repo-readme,repo-pr}.

We use PR encodings of tuples, finite lists, and finite tables. For example,
pairing and its projections may be defined by bounded search within an
elementary numerical bound, and a list may be coded by its length and a
base large enough for its entries. All searches in the finite algorithms below
come with a specified numerical bound. Default values make these operations
total on malformed inputs; no categorical conclusion concerns those defaults.
Arithmetic order, Boolean tests, finite sums and products, finite-list lookup,
and bounded minimization have their familiar PR definitions.

\begin{lemma}[Bounded arithmetic bookkeeping]\label{lem:bounded}
A fixed formula using PR predicates, Boolean connectives and bounded numerical
quantifiers has a PR characteristic function. Bounded scans, table filtering,
and finite-list updates with PR length bounds are PR. The loop invariants used
for these constructions can be expressed quantifier-free and proved by their
indicated base and step calculations.
\end{lemma}
\begin{proof}
Use $1$ for true and $0$ for false. Conjunction is multiplication, negation is
$1-b$, and bounded universal quantification is an iterated product of bits;
bounded existential quantification is one minus the product of complementary
bits. All are primitive recursions. A bounded least witness is obtained by a
scan carrying a flag recording whether an earlier witness was found and its
index. A table update is a scan rebuilding its coded entries.
The defining equations prove the base and step equations for these scans.
For example, the assertion that a returned index is the first successful
one is expressed by its range, its test bit, and a product of failure bits
at earlier indices. This assertion is quantifier-free after replacing the
bounded product by its PR symbol. The same device expresses every bounded
table invariant used below. Consequently induction is applied to a numerical
equation or a Boolean combination of such equations, never to an assertion
about all arbitrary functions. The correspondence between bounded quantifier
notation and its Boolean test follows from the same scan equations, with
any currently discussed index treated as an additional numerical parameter.
\end{proof}

\subsection{Categories as fixed arithmetic presentations}
\begin{definition}[Presented category]\label{def:category}
A presentation $C$ consists of fixed arithmetic formulas
\[
 O_C(a),\qquad H_C(a,b,f),\qquad E_C(a,b,f,g)
\]
and fixed PR terms $\id_C(a)$ and $\operatorname{cmp}_C(a,b,c,g,f)$.
The formula $H_C(a,b,f)$ means that $f$ is an arrow from $a$ to $b$.
The formula $E_C(a,b,f,g)$ compares parallel arrow codes, and is not
necessarily numerical equality or a decidable relation. We write $f\sim g$
for this relation, or simply $f=g$ in categorical equations with the hom
type understood. Write $gf$ or $g\circ f$ for $\operatorname{cmp}(a,b,c,g,f)$.

The presentation comes with iPRA proofs, or explicitly retained hypotheses,
of the following laws. Hom membership entails object membership. On every
hom domain $E_C$ is reflexive, symmetric and transitive. Identity and
composition are well typed. Composition respects $E_C$ in both arguments.
For composable arrows,
\[
 h(gf)\sim(hg)f,\qquad \id_b f\sim f,\qquad f\id_a\sim f.       \tag{C}
\]
All these laws are universally closed implications with the indicated
typing premises. We also require equality judgments to concern well-typed
parallel arrows, or include those premises explicitly whenever they are used.
\end{definition}

The definitions are abbreviations in arithmetic, not new set-existence
axioms. Objects have literal numerical codes; isomorphic objects need not
have equal codes. Only arrows are identified by their hom equivalence.
Different codes for the same informal structure may therefore coexist.
An empty category is allowed, since $O_C$ may be false everywhere.

There are three distinct uses of this definition.
\begin{enumerate}
\item A \emph{fixed presentation}: its formulas and PR terms are chosen
before an arithmetic sentence is formed. Theorems about it are ordinary
iPRA sentences.
\item A \emph{uniform numerical family}: the same formulas and terms have
an extra parameter $p$. The family can be quantified over $p$ because
evaluation of those fixed terms is already PR. Finite table categories
are a particularly useful family of this kind.
\item An \emph{external collection of presentations}: Lean or the
mathematician ranges over formula syntax and PR term syntax. A theorem
constructs an iPRA proof for each instance. This is a proof schema, not
an internal universe whose elements evaluate all PR programs.
\end{enumerate}
Whenever the report says ``every presented category'', the third meaning
is intended unless a numerical family is explicitly supplied.

\begin{theorem}[Specialization of structural proofs]\label{thm:export}
Fix presentations, functors, families and selected universal operations as
defined in this report. A finite typed proof obtained by intuitionistic
logical rules, the supplied laws, and substitution of their PR terms becomes
an iPRA derivation after expanding the abbreviations. If its laws have iPRA
proofs, the conclusion has an iPRA proof; if they are hypotheses, the
conclusion is a relative iPRA theorem with those hypotheses.
No additional induction principle is needed for such a proof.
\end{theorem}
\begin{proof}
Replace every object variable by a number variable with guard $O_C$ and
every arrow variable by a number variable with its hom guard. Replace
categorical equality by $E_C$, and categorical operations by their fixed
terms. A well-typed use of reflexivity, symmetry, transitivity, congruence,
associativity or identity is exactly an instance of a supplied law.
Composition of proof steps is implication elimination or cut. An explicit
witness term justifies existential introduction; an existential hypothesis
may be eliminated locally with the usual eigenvariable condition.
Generalization introduces only numerical quantifiers and their guards.
Thus each line translates to an allowed inference from earlier lines.
There are finitely many lines in a fixed proof, so this construction needs
no induction in the object theory. Formally, induction on proof syntax is
performed in the metatheory to describe the translator, not admitted as
an extra object-language rule. Every induction in a concrete arithmetic
implementation must separately satisfy (I).
\end{proof}

This theorem is intentionally a specialization theorem. It does not assert
that arbitrary Lean theorems are iPRA theorems. In particular, a Lean field
$\N\to\N$ does not certify a PR implementation, and a proof using induction
on a proposition containing unbounded quantifiers does not pass this audit
merely because its conclusion is arithmetical. ProveIt's
\texttt{SetoidInterpretation.sound} gives a more general syntactic
interpretation mechanism when all domain and equality compatibility laws
are supplied \cite{repo-setoid}. The present proofs use the smaller explicit
typed calculus just described.

\subsection{Examples that do not reduce the subject to finite categories}
Every fixed PR monoid $(M,\cdot,e)$ gives a one-object presentation: arrows
are its coded elements, identity is $e$, composition is multiplication.
The additive monoid of natural numbers is the elementary example already
instantiated in \texttt{PresentedCategory.natAdd\_laws}.
A PR preorder gives a thin category: take $H(a,b,f)$ to assert the order
relation and $f=0$; all parallel arrows are equal.
Integer matrices give an infinite presentation with dimension codes as
objects, finite coded matrices as arrows and matrix multiplication as
composition. Its associativity reduces to finite distributivity and
reindexing of bounded sums, hence to Lemma~\ref{lem:bounded} and the PR
ring laws. This example has infinitely many objects and arrows.

Equality may also have unbounded logical complexity. For instance,
the equality of two represented numerical sequences may be
$\forall n\,a(n)=b(n)$, and regular real names may be compared by
an arithmetical approximation condition. For a fixed family of PR names,
operations and a modulus, these are legitimate formulas. Establishing
the analytic category laws is a separate mathematical task; the structural
theorems do not decide equality or assert those laws without proof.

\subsection{Relational programs: a useful extension with a different interface}
There is a further way to have numerical program codes without a total
PR evaluator. Fix a deterministic register-machine language and a PR
predicate $T(e,x,y,t)$ saying that $t$ is a complete halting computation
of code $e$ on input $x$ with output $y$. Its finite-transition check is
PR. Let
\[
 \operatorname{Tot}(e)\equiv\forall x\,\exists y,t\,T(e,x,y,t).
\]
Take one object, arrows $e$ satisfying $\operatorname{Tot}(e)$, and
\[
 e\sim e'\quad\Longleftrightarrow\quad
 \forall x,y,z,t,u\,[T(e,x,y,t)\wedge T(e',x,z,u)\to y=z].       \tag{R}
\]
Identity and composition are PR \emph{code constructors}. To make the
presentation fully explicit, composition first runs its first component,
then runs its second component on the result, with disjoint renamed
work registers. Its trace certificate stores the two component traces
and their joining configuration. Inlining these tagged certificates into
ordinary transition traces is a PR operation.

\begin{proposition}[Relational composition]
Given the finite trace-composition and decomposition lemmas for this fixed
machine language, (R) defines a presented category. The proof uses neither
choice nor an internal assertion that all program codes terminate.
\end{proposition}
\begin{proof}
Determinism gives reflexivity of (R), and symmetry is immediate. For
transitivity on total arrows, obtain an output trace of the middle program
at the particular input $x$; the two equality hypotheses compare the first
and last outputs with that middle output. This is one existential
elimination, not a simultaneous selection over all inputs.
For totality of a composite, fix $x$, use totality of $f$ to obtain $y,t$,
then totality of $g$ at that $y$ to obtain $z,u$, and concatenate the two
traces. Equality is preserved by composition by applying (R) first to
the input-side results and then to the output-side results. Both
associations of three programs run the same three computations; decomposing
their traces and using determinism proves equality. Identity traces give
the identity laws. Each argument compares only finitely many runs at a
time. The trace lemmas themselves are proved by quantifier-free induction
on coded finite traces and their bounded local checks.
\end{proof}

This optional interface is a precise conditional construction: the report
does not ship a newly verified register-machine trace library. It explains
why the evaluator obstruction below does not rule out arithmetically
defined collections of total programs. It also explains the limitation:
$\operatorname{Tot}(e)$ is not a PR test, and extracting a table of outputs
for a variable number of such runs requires its own trace, clock, or
batch-totality proof. The structural development uses fixed PR operations
on arrow codes, not a claim that semantic application is itself PR.

\subsection{Elementary consequences and categorical constructions}
An isomorphism $f:a\to b$ consists of a specified $g:b\to a$ with
$gf=\id_a$ and $fg=\id_b$. A monomorphism cancels on the left and an
epimorphism cancels on the right; these are guarded arithmetic formulas.

\begin{proposition}\label{prop:elementary}
In every presented category inverses are unique. Every section is monic,
every retraction is epic, and every monic retraction is an isomorphism.
Functors preserve isomorphisms; faithful functors reflect equality of
parallel arrows.
\end{proposition}
\begin{proof}
If $gf=\id_a$ and $fh=\id_b$, then
$g=g\id_b=g(fh)=(gf)h=\id_a h=h$.
If $rs=\id_a$ and $su=sv$, then
$u=(rs)u=r(su)=r(sv)=(rs)v=v$.
Reversing arrows gives the retraction assertion.
If $fr=\id_b$ and $f$ is monic, then
$f(rf)=(fr)f=f=f\id_a$, so $rf=\id_a$.
Applying a functor to the two inverse equations gives the inverse equations
for its image. Reflection of equality is exactly faithfulness.
\end{proof}
Being both monic and epic does not in general imply invertibility: in the
thin category on $0<1$, the arrow $0\to1$ cancels on both sides, but there
is no reverse arrow. This finite counterexample is already available in iPRA.

The opposite category reverses hom domains and composition, leaving all
codes unchanged. Its laws are the reversed laws of $C$. Reversing a fixed
proof line by line therefore proves the dual theorem. The product of two
categories pairs object and arrow codes and acts componentwise.
For a fixed object $b$, the slice $C/b$ has objects $(a,u:a\to b)$ and
arrows $f:(a,u)\to(a',u')$ satisfying $u'f=u$.
The coslice is dual. The comma category $(F\downarrow G)$ has objects
$(a,b,u:Fa\to Gb)$ and arrows $(f:a\to a',g:b\to b')$ satisfying
$Gg\,u=u'\,Ff$. Equality is componentwise.
For identity the square is an identity equation. For composition, if
$Gg\,u=u'Ff$ and $Gg'\,u'=u''Ff'$, then
$G(g'g)u=Gg'(Gg\,u)=Gg'(u'Ff)=u''F(f'f)$.
Associativity and both unit laws are componentwise laws. Thus all these
constructions have PR code operations and arithmetic domain formulas;
no collection of arbitrary sets or functions was formed.
